\documentclass[11pt]{article}
\usepackage{mathunal-base}
\mathunalfuente{serif}

\renewcommand{\examtitulo}{Tercer Examen Parcial de Métodos Numéricos (25\%)}
\renewcommand{\examfecha}{24 de noviembre de 2007}

\begin{document}

\examheader

\vspace{4pt}
\noindent Nombre: \dotfill\ Carné: \dotfill\ Grupo: \dotfill

\vspace{8pt}
\noindent\textbf{Duración: 1 hora y 50 minutos}

\begin{center}\textbf{Toda respuesta debe estar debidamente justificada}\end{center}

\begin{enumerate}
\item{} (20\%) Considere el P.V.I.
\[
\left\{
\begin{aligned}
y' &= \frac{4t^3y}{1+t^4}, \qquad 1\leq t\leq 3\\
y(1) &= 2
\end{aligned}
\right.
\]
\begin{enumerate}
\item{} (15\%) ¿Satisface la función $f(t,y)=\dfrac{4t^3y}{1+t^4}$ una condición de Lipschitz con respecto a la variable $y$ en la región $R=\{(t,y)\ /\ 1\leq t\leq 3,\ -\infty<y<\infty\}$? En caso afirmativo, determine el menor valor de la constante de Lipschitz.

\vspace{4cm}

\item{} (5\%) ¿Tiene este problema con valor inicial una única solución?

\vspace{2cm}
\end{enumerate}

\item{} (15\%) Considere el P.V.I.
\[
\left\{
\begin{aligned}
y' &= te^{-2y}, \qquad t\geq 0\\
y(0) &= 1
\end{aligned}
\right.
\]
\noindent Para \textbf{este} problema de valor inicial, escriba la fórmula de avance del método de Taylor de orden 2.

\vspace{5cm}

\item{} Considere el problema de contorno o de valor en la frontera
\[
\left\{
\begin{aligned}
y''-e^{-t}y'-y &= t^2, \qquad 2\leq t\leq 4\\
y(2) &= -1,\quad y(4)=3.5
\end{aligned}
\right.
\]
\begin{enumerate}
\item{} (10\%) ¿Tiene este problema de contorno o de valores en la frontera una única solución?

\vspace{3cm}

\item{} (20\%) Unas soluciones aproximadas de los problemas de valor inicial asociados:
\[
\left\{
\begin{aligned}
u''-e^{-t}u'-u &= t^2\\
u(2)&=-1\\
u'(2)&=0
\end{aligned}
\right.
\qquad 2\leq t\leq 4
\qquad
\left\{
\begin{aligned}
w''-e^{-t}w'-w &= 0\\
w(2)&=0\\
w'(2)&=1
\end{aligned}
\right.
\qquad 2\leq t\leq 4
\]

\vspace{4pt}
\noindent con tamaño de paso $h=0.5$ empleando el método de Heun están dadas en las siguientes tablas.

\vspace{6pt}
\begin{center}
\begin{tabular}{|c|c|c|c|c|c|}
\hline
 & $t=2$ & $t=2.5$ & $t=3$ & $t=3.5$ & $t=4$ \\ \hline
$u(t)$ & $-1$ & $-0.6250$ & $1.1462$ & $5.4316$ & $14.0646$ \\ \hline
$u'(t)$ & $0$ & $2.0933$ & $5.9600$ & $\mathbf{D}$ & $24.3172$ \\ \hline
\end{tabular}

\vspace{6pt}
\begin{tabular}{|c|c|c|c|c|c|}
\hline
 & $t=2$ & $t=2.5$ & $t=3$ & $t=3.5$ & $t=4$ \\ \hline
$w(t)$ & $0$ & $0.5169$ & $1.1840$ & $2.1569$ & $3.6671$ \\ \hline
$w'(t)$ & $1$ & $1.1807$ & $\mathbf{C}$ & $2.4626$ & $3.8839$ \\ \hline
\end{tabular}
\end{center}

\vspace{4pt}
\noindent Hallar los valores de $C$ y $D$ \textbf{definiendo claramente} las variables y funciones empleadas.

\vspace{6cm}

\item{} (10\%) \textbf{Hallar} una aproximación a la solución del problema de contorno o de valores en la frontera con el mismo tamaño de paso.

\vspace{5cm}
\end{enumerate}

\item{} (20\%) Considere el problema
\[
\left\{
\begin{aligned}
\frac{\partial u}{\partial t} - \frac{1}{9}\frac{\partial^2 u}{\partial x^2} &= 0, \qquad 0<x<1,\ t>0\\
u(0,t) &= e^t, \qquad t\geq 0\\
u(1,t) &= t^2+1, \qquad t\geq 0\\
u(x,0) &= \cos(2\pi x), \qquad 0\leq x\leq 1
\end{aligned}
\right.
\]
\noindent Empleando el método de diferencias finitas progresivas con $h=\frac{1}{5}$ y $k=0.15$, hallar aproximaciones de $u(x,t)$ para:
\begin{enumerate}
\item{} (4\%) $u\left(\frac{3}{5},0\right)$.
\item{} (4\%) $u(1,0.3)$.
\item{} (12\%) $u\left(\frac{1}{5},0.15\right)$.
\end{enumerate}

\vspace{6cm}
\end{enumerate}

\examfooter
\end{document}
